\chapter{Implementierung}

Die technische Umsetzung der Webanwendung orientiert sich direkt an den im vorangegangenen Kapitel definierten fachlichen und didaktischen Anforderungen. Ziel dieses Kapitels ist die konkrete Realisierung eines interaktiven Scheduling-Visualizers, der fachlich korrekte Simulationsergebnisse mit einer verständlichen und steuerbaren Oberfläche verbindet. Dabei bleibt die fachliche Logik von der Präsentation getrennt, sodass die Visualisierung immer auf nachvollziehbare Berechnungsergebnisse zurückgeführt werden kann.

\section{Technische Grundlagen und verwendete Technologien}

Die Anwendung wurde als clientseitige Webanwendung mit Vue 3 und TypeScript umgesetzt. Vue 3 dient als Basis für die komponentenbasierte UI-Struktur, während TypeScript die Modellierung von Tasks, Events, Snapshots, Runs und Kennzahlen typensicher macht. Vite übernimmt den Build- und Entwicklungsprozess, Vitest dient der automatisierten Verifikation der Kernlogik und GSAP ergänzt die Oberfläche mit gezielten Übergängen und einer klaren dynamischen Präsentation \, ohne die fachliche Logik der Simulation zu beeinflussen.

Die zentrale technische Entscheidung ist die Trennung von fachlicher Modellierung, Laufsteuerung und Visualisierung. Dadurch entsteht ein klarer Datenfluss: Eingabedaten werden in eine deterministische Simulation überführt, die Ergebnisse werden als Ereignisliste, Segmente und Kennzahlen modelliert und anschließend in der Oberfläche als Gantt-Darstellung, Queue-Ansicht und Vergleichsübersicht visualisiert.

\section{Architektur der Anwendung}

Die Architektur folgt dem in Kapitel 4 beschriebenen Schichtenmodell mit drei Ebenen:

\begin{itemize}
    \item \textbf{Szenario- und Konfigurationslogik:} Verwaltung von Taskdaten, Parametern, Szenario-Definitionen und gespeicherten Laufkonfigurationen.
    \item \textbf{Simulationskern:} Deterministische Berechnung der Scheduling-Strategien auf Basis eines gemeinsamen Zustandsmodells.
    \item \textbf{Darstellung und Steuerung:} Gantt-Darstellung, Warteschlangenvisualisierung, Kennzahlbereiche und Interaktionssteuerung für Abspielen, Pausieren und Vergleichen.
\end{itemize}

Diese Gliederung ist nicht nur aus Sicht des Software-Engineering sinnvoll, sondern entspricht auch den didaktischen Anforderungen der Anwendung. Der Nutzer kann den fachlichen Ablauf nachvollziehen, ohne die technische Struktur der Oberfläche direkt mit der Berechnung vermischen zu müssen.

\section{Datenmodell und Simulationskern}

Die fachliche Grundlage bildet ein gemeinsames Task-Modell mit Attributen wie Ankunftszeit, Burst-Zeit, Priorität, Verlauf, Restzeit und Status. Auf dieser Grundlage werden die einzelnen Scheduling-Verfahren als Varianten eines gemeinsamen Ablaufmodells umgesetzt. Dadurch bleiben die Verfahren unter identischen Bedingungen vergleichbar, und die Unterschiede im Ergebnis sind fachlich auf die jeweilige Algorithmik zurückführbar.

Der Kern erzeugt nicht nur ein Endergebnis, sondern mehrere, logisch zusammengehörige Datensätze:

\begin{itemize}
    \item \textbf{Ereignislog:} chronologische Information über Dispatches, Präemptionen, Kontextwechsel und Abschluss.
    \item \textbf{Segmente:} zeitlich abgegrenzte Ausführungsschritte für die Gantt-Darstellung.
    \item \textbf{Snapshots:} Zustandsbilder für die Wiedergabe und die genaue Zeitmarkierung im Ablauf.
    \item \textbf{Metriken:} aggregierte Kennzahlen für Wartezeit, Durchlaufzeit, Reaktionszeit und Kontextwechsel.
\end{itemize}

Diese Struktur ist besonders wichtig für die fachliche Nachvollziehbarkeit: Die Visualisierung bezieht sich immer auf dieselben Daten, die auch der kennzahlenbasierte Vergleich verwendet. So bleiben die Darstellung und die Bewertung konsistent.

\section{Umsetzung der Scheduling-Verfahren}

Die Simulation modelliert das Verhalten der wichtigsten präemptiven Verfahren innerhalb eines identischen Simulationsrahmens. Round Robin wird durch Zeitquantum-Logik realisiert, Strict Priority durch prioritätsbasierte Auswahl mit Präemptionsentscheidungen, LCFS über LIFO-basierte Warteschlangenlogik und MLFQ über mehrere Ebenen mit reduzierter Priorität bei längeren Ausführungsintervallen.

Die zentrale Gemeinsamkeit aller Verfahren liegt in der konsistenten Tick- und Ereignislogik. Jeder Lauf wird auf Grundlage desselben Szenarios und derselben Zustandsbasis berechnet, sodass Unterschiede in Laufzeit, Wartezeit und Durchlaufzeit tatsächlich auf die Strategiedefinitionen zurückzuführen sind und nicht auf abweichende Berechnungsmethoden.

\section{Umsetzung der Visualisierung}

Die Oberfläche ist in mehrere, logische Bereiche aufgeteilt: Szenario- und Parameterwahl, Laufsteuerung, Fokusansicht, Vergleichsübersicht und Kennzahlenbereich. Die Gantt-Zeitachse bildet den Ablauf eines einzelnen Laufes ab, die Stack-Visualisierung zeigt den aktuellen Zustand der Warteschlange und die Kennzahlen verknüpfen den Lauf mit relevanten fachlichen Metriken.

Die Anwendung nutzt SVG-basierte Darstellungen, um Zeitabschnitte, Segmentgrenzen, Ankunftsereignisse und Präemptionen präzise in der Oberfläche zu markieren. Dadurch entsteht eine visuell klare, fachlich genaue Darstellung, die zugleich als Lehrmittel genutzt werden kann. Die Visualisierung ist daher nicht nur dekorativ, sondern trägt unmittelbar zum Verständnis von Scheduling-Entscheidungen bei.

\section{Interaktionsmechanismen und Nutzerworkflow}

Die Interaktion folgt einem didaktisch klaren Ablauf: Szenario definieren, Algorithmus wählen, Lauf steuern, Ergebnis interpretieren und vergleichen. Die Wiedergabelogik wird getrennt von der Simulation modelliert. Sie liest vorberechnete Snapshots und zeigt nur den jeweils aktuellen Ablaufzustand an, ohne den fachlichen Berechnungsprozess erneut zu starten.

Wichtige Interaktionsfunktionen sind Start, Pause, Schrittmodus, Geschwindigkeitsanpassung und Zeitsprung. Ergänzend gibt es einen Vergleichsmodus, der mehrere Läufe auf identischer Eingabebasis nebeneinander darstellt. Dadurch wird die fachliche Bedeutung der Algorithmuswahl unmittelbar erfahrbar und mit den entsprechenden Kennzahlen verknüpft.

\section{Besondere Designentscheidungen}

Eine zentrale Designentscheidung war die Verbindung von fachlicher Korrektheit und didaktischer Transparenz. Die Anwendung zeigt nicht nur das Ergebnis, sondern auch die Zustandswechsel, die zu einer Entscheidung geführt haben. Dadurch wird aus einer bloßen Simulation ein erklärbares Lehrwerkzeug.

Ein weiterer wichtiger Punkt ist die klare Trennung zwischen Berechnung und Vergleich. Die Laufberechnung erzeugt die Basisdaten, während die Vergleichsansicht diese Daten auswertet und visuell gegenüberstellt. Diese Struktur verbessert die Wiederverwendbarkeit, vereinfacht die Testbarkeit und erleichtert zukünftige Erweiterungen um weitere Algorithmen oder Lernpfade.
